\documentclass[14 pt]{extarticle}

	\usepackage[frenchb]{babel}
	\usepackage[utf8]{inputenc}  
	\usepackage[T1]{fontenc}
	\usepackage{amssymb}
	\usepackage[mathscr]{euscript}
	\usepackage{stmaryrd}
	\usepackage{amsmath}
	\usepackage{tikz}
	\usepackage[all,cmtip]{xy}
	\usepackage{amsthm}
	\usepackage{varioref}
	\usepackage{geometry}
	\geometry{a4paper}
	\usepackage{lmodern}
	\usepackage{hyperref}
	\usepackage{array}
	 \usepackage{fancyhdr}
\renewcommand{\theenumi}{\alph{enumi})}
	\pagestyle{fancy}
	\theoremstyle{plain}
	\fancyfoot[C]{} 
	\fancyhead[L]{Contrôle}
	\fancyhead[R]{12 juin 2025}\geometry{
 a4paper,
 total={170mm,257mm},
 left=20mm,
 top=20mm,
 }
	
	
	\title{Contrôle proportionnalité}
	\date{}
	\begin{document}
\emph{Toute réponse doit être soigneusement justifiée.}
 \subsection*{Exercice 1 - 5 points}
 Dans une recette pour $4$ personnes, on utilise $300$ g de farine, $250$ g de sucre, $4$ oeufs et $6$ g de levure. 
 Donner les quantités pour $6$ personnes. 
 \subsection*{Exercice 2 - 5 points}
 Remplir les tableaux de proportionnalité suivants : 
 \begin{enumerate}
 \item 
 $\renewcommand{\arraystretch}{1}
\begin{array}{|c|c|c|c|}
\hline
\phantom{000}10\phantom{000} & \phantom{000}6\phantom{000} & \phantom{000}4\phantom{000}& \phantom{000}16\phantom{000}\\
\hline
\phantom{000}\phantom{000} & \phantom{000}57\phantom{000} & \phantom{000}38\phantom{000}& \phantom{000}\phantom{000}\\
\hline
\end{array}
\renewcommand{\arraystretch}{1}$
\item $\renewcommand{\arraystretch}{1}
\begin{array}{|c|c|c|}
\hline
\phantom{000}3\phantom{000} & \phantom{000}\phantom{000} & \phantom{000}10,5\phantom{000}\\
\hline
4 & 120 & \\
\hline
\end{array}
\renewcommand{\arraystretch}{1}$
 \end{enumerate}
 \subsection*{Exercice 3 - 5 points}
 Calculer : \begin{enumerate}
 \item $30\%$ de $42$ euros. 
 \item $15\%$ de $54$ grammes. 
 \end{enumerate}
 Parmi les 26 lettres de l'alphabet $6$ sont des voyelles. Quel pourcentage de l'alphabet représentent donc les voyelles ?  
 \subsection*{Exercice 4 - 5 points}
Dans une classe, il y  a 12 filles et 18 garçons.
\begin{enumerate}
\item  Quel est le pourcentage de filles ?
\item On sait que cette classe contient $15$ \% des élèves du collège. Combien y a-t-il d'élèves ? 
\end{enumerate}

 	\end{document}
 	\newpage
